\section{Limitations}
\label{sec:limits}

The mesh is a far-field method. Its limits are set by $r/h$ and $\gamma$,
and they are measured, not assumed. The mesh alone does not reach $10^{-3}$
where the slices cut the bunch core (Section~\ref{sec:slices}); that is the
single-mesh limitation the slices display. The opt-in near shell
(Section~\ref{sec:nearsheet}) assembles exactly to roundoff, so at the
measured radii the residual is the mesh outside the shell, not the tile
subtraction. The perpendicular direction at $\gamma=10$ does not converge at
the natural cell size: $1.5\times10^{-2}$ at $d=48$, aspect 1, against the
point field, of which the TSC smoothing alone is $1.48\times10^{-2}$
(Table~\ref{tab:d48}) \citep{N2b,D11}. At $\gamma=100$ the perpendicular error
is $3.75\times10^{-1}$ at $d=48$ and does not converge as $d$ grows
\citep{N2b,D8}. At aspect 4, $\gamma=1$ already fails the gate at $d=48$
($2.4\times10^{-3}$) \citep{N2b}. The measured near-shell cost ladder does
not fit the $2\times10^{6}$ / $10^{-3}$ / $300$\,s envelope at production
size (Section~\ref{sec:cost}). The detailed discussion is in
Section~\ref{sec:discussion}.
